\documentclass[12pt,a4paper]{article}
\usepackage[utf8]{inputenc}
\usepackage[spanish]{babel}
\usepackage[margin=2.5cm]{geometry}
\usepackage{graphicx}
\usepackage{booktabs}
\usepackage{hyperref}
\usepackage[backend=biber,style=ieee]{biblatex}
\addbibresource{referencias.bib}

\begin{document}

\begin{titlepage}
    \centering
    {\Large \textbf{UNIVERSIDAD TÉCNICA ESTATAL DE QUEVEDO}}\\[0.3cm]
    {\large Facultad de Ciencias de la Ingeniería}\\[0.2cm]
    {\large Carrera de Ingeniería en Software / Telemática}\\[1.5cm]
    {\large \textbf{Asignatura:} Sistemas Distribuidos}\\[0.3cm]
    {\large \textbf{Actividad:} TA-IND-01 — Informe Académico Individual}\\[1.5cm]
    {\LARGE \textbf{Análisis crítico comparativo de Apache Kafka y RabbitMQ}}\\[2cm]
    \begin{flushleft}
        \textbf{Estudiante:} [Manotoa Alava Carlos Alexander]\\
        \textbf{Fecha:} 15 de noviembre de 2026\\
        \textbf{Repositorio:} \url{https://github.com/BelindaTZ/TA-IND-01-Informe-Acad-mico-Individual---CGMT.git}
    \end{flushleft}
\end{titlepage}

\section*{Resumen}
% De 150 a 200 palabras

\section{Introducción}
\section{Arquitectura del sistema}
\section{Análisis con los conceptos de la Unidad 1}
\subsection{Transparencias de distribución}
\subsection{Teorema CAP y extensión PACELC}
\subsection{Modelos de comunicación}
\subsection{Tolerancia a fallos}
\section{Experimentación propia}
\section{Ejemplos del entorno del estudiante}
\section{Análisis comparativo}
\section{Postura crítica}
\section{Conclusiones}
\section{Declaración de uso de inteligencia artificial}

\printbibliography[title={Referencias}]

\appendix
\section{Evolución del esquema de arquitectura}
\section{Evidencias de lectura de fuentes}

\end{document}